\section*{Capítulo \ref{ch:b3:group-theory-applied} --- Teoría de grupos aplicada}

\begin{solution}{exo:b3:group-theory-applied:1}
$n_{A_1} = \frac14(5 - 1 + 3 + 1) = 2$, $n_{A_2} = \frac14(5 - 1 - 3 - 1) = 0$,
$n_{B_1} = \frac14(5 + 1 + 3 - 1) = 2$, $n_{B_2} = \frac14(5 + 1 - 3 + 1) = 1$:
$\Gamma = 2A_1 + 2B_1 + B_2$, de dimensión 5.
\end{solution}

\begin{solution}{exo:b3:group-theory-applied:2}
El mismo recuento de átomos que quedan en su sitio que para el agua: $\Gamma_{\mathrm{vib}} = 2A_1 + B_1$. $A_1$ ($z$) y
$B_1$ ($x$) son \omterm{def:b3:group-theory-applied:activity}{activos en IR}, y ambos son también \omterm{def:b3:group-theory-applied:activity}{activos en Raman} ($x^2$, $xz$): tres
bandas en cada espectro.
\end{solution}

\begin{solution}{exo:b3:group-theory-applied:3}
$B_1\otimes B_2$: \omterm{def:b3:point-groups:representation}{caracteres} $(1, 1, -1, -1)$ $= A_2$. $A_2\otimes B_1$: $(1, -1, -1,
1) = B_2$. En $C_{3v}$, $E\otimes E$ tiene los \omterm{def:b3:point-groups:representation}{caracteres} $(4, 1, 0)$; reducción: $A_1 + A_2 +
E$.
\end{solution}

\begin{solution}{exo:b3:group-theory-applied:4}
Infrarrojo: los dos modos $t_2$, 3019 y \qty{1306}{cm^{-1}}. Raman: los cuatro ($a_1$,
$e$ y $t_2$ se transforman todos como funciones cuadráticas). No hay \omterm{def:b3:point-groups:inversion}{centro de inversión}, de modo que
las coincidencias están permitidas.
\end{solution}

\begin{solution}{exo:b3:group-theory-applied:5}
Átomos que quedan en su sitio, $(4, 1, 2, 4, 1, 2)$, multiplicados por $\chi_{xyz} = (3, 0, -1, 1, -2, 1)$, dan
$\chi_{3N} = (12, 0, -2, 4, -2, 2)$, que se reduce a $A_1' + A_2' + 3E' + 2A_2'' +
E''$. Se restan $E' + A_2''$ (traslaciones) y $A_2' + E''$ (rotaciones): $A_1' + 2E'
+ A_2''$. IR: $2E' + A_2''$, tres bandas; Raman: $A_1' + 2E'$, tres bandas. $A_2''$
(flexión fuera del plano) solo es \omterm{def:b3:group-theory-applied:activity}{activo en IR}, $A_1'$ (tensión simétrica) solo en Raman, y los modos
$E'$ aparecen en ambos.
\end{solution}

\begin{solution}{exo:b3:group-theory-applied:6}
IR: los dos modos $T_{1u}$, dos bandas. Raman: $A_{1g}$, $E_g$, $T_{2g}$, tres bandas.
$T_{2u}$ es silencioso. \ce{SF6} tiene \omterm{def:b3:point-groups:inversion}{centro de inversión}: exclusión mutua.
\end{solution}

\begin{solution}{exo:b3:group-theory-applied:7}
Con $E$, $C_3$, $C_3^2$, la función $h_2$ pasa a $h_2$, $h_3$, $h_1$ (\omterm{def:b3:point-groups:representation}{caracteres} 2,
$-1$, $-1$); las reflexiones contribuyen con 0. $\hat P_Eh_2 \propto 2h_2 - h_3 - h_1$.
Su solapamiento con $2h_1 - h_2 - h_3$ (despreciando los solapamientos atómicos) es $-3$, y esta
última tiene norma $\sqrt6$; restando la proyección, $(2h_2 - h_1 - h_3) +
\frac12(2h_1 - h_2 - h_3) = \frac32(h_2 - h_3)$.
\end{solution}

\begin{solution}{exo:b3:group-theory-applied:8}
fac ($C_{3v}$, \omterm{def:b3:point-groups:representation}{caracteres} $3, 0, 1$): $A_1 + E$, dos bandas IR. mer ($C_{2v}$,
\omterm{def:b3:point-groups:representation}{caracteres} $3, 1, 3, 1$): $2A_1 + B_1$, tres bandas. \ce{M(CO)6} ($O_h$): $A_{1g} + E_g
+ T_{1u}$, una banda IR ($T_{1u}$).
\end{solution}

\begin{solution}{exo:b3:group-theory-applied:9}
\omterm{def:b3:point-groups:representation}{Caracteres} $6, 0, 0, 2, 2, 0, 0, 0, 4, 2$: solo $E$, los $C_4$ y $C_2$ según los
ejes (dos ligandos que quedan en su sitio), los tres $\sigma_h$ (cuatro) y los seis $\sigma_d$
(dos) dejan ligandos en su sitio. La reducción da $A_{1g} + E_g + T_{1u}$. Se combinan los orbitales $s$
($a_{1g}$), $d_{z^2}$ y $d_{x^2-y^2}$ ($e_g$), y $p_x$, $p_y$, $p_z$ ($t_{1u}$) del metal;
$t_{2g}$ queda no enlazante.
\end{solution}

\begin{solution}{exo:b3:group-theory-applied:10}
$n_i = \frac1h(1\cdot h\cdot\chi_i(E)) = d_i$. La dimensión de la representación
regular es $h$, y también $\sum_in_id_i = \sum_id_i^2$.
\end{solution}

\begin{solution}{exo:b3:group-theory-applied:11}
En $D_{2h}$ (molécula según $z$), los átomos que quedan en su sitio son
$(4, 4, 0, 0, 0, 0, 4, 4)$, de donde $\chi_{3N}
= (12, -4, 0, 0, 0, 0, 4, 4)$, que se reduce a $2A_g + 2B_{2g} + 2B_{3g} + 2B_{1u} + 2B_{2u}
+ 2B_{3u}$. Se restan las traslaciones $B_{1u} + B_{2u} + B_{3u}$ y las dos rotaciones
$B_{2g} + B_{3g}$: $2A_g + B_{2g} + B_{3g} + B_{1u} + B_{2u} + B_{3u}$, siete modos.
En $D_{\infty h}$: $2\Sigma_g^+ + \Pi_g + \Sigma_u^+ + \Pi_u$. IR: $\Sigma_u^+$ y $\Pi_u$
(dos bandas); Raman: $2\Sigma_g^+$ y $\Pi_g$ (tres bandas); ninguna coincidencia.
\end{solution}

\begin{solution}{exo:b3:group-theory-applied:12}
$1b_2 \to 4a_1$: $B_2$, $y$, permitida, perpendicular al plano. $3a_1 \to 4a_1$:
$A_1$, $z$, permitida, según el eje $C_2$. $1b_2 \to 2b_1$: $A_2$, prohibida.
$1b_1 \to 4a_1$: $B_1$, $x$, permitida, en el plano y perpendicular al eje.
\end{solution}

\begin{solution}{pb:b3:group-theory-applied:1}
\textbf{1.} cis: los dos CO en vértices contiguos; trans: opuestos; fac: tres
CO en una cara; mer: tres CO en un meridiano (dos en trans y uno entre ellos).
\textbf{2.} $C_{2v}$.
\textbf{3.} $D_{4h}$.
\textbf{4.} $C_{3v}$.
\textbf{5.} $C_{2v}$.
\textbf{6.} Solo trans-\ce{ML4(CO)2}.
\textbf{7.} Una operación que lleva un CO sobre otro saca su vector de la
diagonal (contribución 0); una que deja un CO en su sitio transforma su vector de enlace
en sí mismo (+1); ninguna invierte un vector C--O.
\textbf{8.} \omterm{def:b3:point-groups:representation}{Caracteres} $(2, 0, 2, 0)$, con $\sigma_v(xz)$ el plano de los dos CO: $A_1 +
B_1$.
\textbf{9.} \omterm{def:b3:point-groups:representation}{Caracteres} 2 para $E$, $C_4$, $C_2$, $\sigma_v$, $\sigma_d$, y 0 para las demás:
$A_{1g} + A_{2u}$.
\textbf{10.} $(3, 0, 1)$: $A_1 + E$.
\textbf{11.} $(3, 1, 3, 1)$: $2A_1 + B_1$.
\textbf{12.} 2, 2, 3, 3: una dimensión por CO.
\textbf{13.} \omterm{def:b3:group-theory-applied:activity}{Activas en IR}: cis 2, trans 1, fac 2, mer 3.
\textbf{14.} \omterm{def:b3:group-theory-applied:activity}{Activas en Raman}: cis 2, trans 1, fac 2, mer 3.
\textbf{15.} El trans: su banda IR ($A_{2u}$) y su banda Raman ($A_{1g}$) son modos
distintos.
\textbf{16.} $A_{2u}$ cambia de signo con $i$: un CO se alarga mientras el otro
se acorta, en oposición de fase.
\textbf{17.} $E$ es un par degenerado: dos modos de la misma frecuencia dan una sola
banda, más la banda $A_1$.
\textbf{18.} Un isómero trans-disustituido y centrosimétrico.
\textbf{19.} Tres bandas IR: mer (fac daría dos).
\textbf{20.} Sí: mer tiene tres modos CO \omterm{def:b3:group-theory-applied:activity}{activos en Raman} a las mismas frecuencias,
por no ser centrosimétrico; fac presentaría dos.
\textbf{21.} $A_1$: la tensión en fase es totalmente simétrica.
\textbf{22.} Una mezcla fac/mer presentaría hasta cinco bandas IR, con las bandas
de fac en posiciones distintas; tres bandas limpias corresponden a un solo isómero. Un
espectro de RMN de ${}^{13}$C lo decidiría: dos señales de carbonilo en razón 2:1 para
mer, una sola señal para fac.
\textbf{23.} \textbf{El isómero mer tiene tres tensiones CO \omterm{def:b3:group-theory-applied:activity}{activas en IR}, y el isómero
fac dos}: el producto es mer.
\end{solution}
